% !TeX program = xelatex
% UTF-8 transcription; source readings and arrows are retained, not medically corrected.
% pic01.jpg = 20261013_023756737.jpg
% pic02.jpg = 20261013_023811226.jpg
% pic03.jpg = 20261013_023815907.jpg
% pic04.jpg = 20261013_023825196.jpg
% pic05.jpg = 20261013_023831509.jpg
% pic06.jpg = 20261013_023838225.jpg
% pic07.jpg = 20261013_023842689.jpg
% pic08.jpg = 20261013_023847687.jpg
% Compile standalone with XeLaTeX, or input into an existing XePersian document.
\ifdefined\settextfont
  \def\AdrenalStandalone{0}
\else
  \def\AdrenalStandalone{1}
  \documentclass[12pt,a4paper]{article}
  \usepackage[margin=2cm]{geometry}
  \usepackage{amsmath,amssymb}
  \usepackage{array,longtable}
  \usepackage{xcolor}
  \usepackage{xepersian}
  \IfFontExistsTF{Amiri}{\settextfont{Amiri}}{\settextfont{FreeSerif}}
  \setlatintextfont{FreeSerif}
\fi
\providecommand{\en}[1]{\lr{#1}}
\newcommand{\admark}[1]{\textbf{#1}}
\newcommand{\aduncertain}[1]{#1\,\textsuperscript{[خوانش نامطمئن]}}
\newcommand{\adarrow}{\en{$\longleftarrow$}}
\newcommand{\adpage}[1]{\clearpage\section*{آدرنال ـ صفحهٔ #1}}
\newcommand{\adprinted}{\noindent{\footnotesize\en{Date:}\quad\en{Subject:}\par}}
\newcommand{\adfooter}[1]{\par\smallskip\noindent{\footnotesize\en{GNOTE}\quad #1\par}}
\ifnum\AdrenalStandalone=1 \begin{document}\fi

\adpage{۱}
\adprinted
\begin{itemize}
\item شایع‌ترین علت نارسایی \admark{اولیه} آدرنال \adarrow\ \admark{بیماری آدیسون}.
\item آدیسون \adarrow\ گلوکوکورتیکوئید \en{$\downarrow$}، مینرالوکورتیکوئید \en{$\downarrow$}:
  \begin{itemize}
  \item \en{$\mathrm{Na}\downarrow$}، \en{$\mathrm{K}\uparrow\uparrow$}.
  \item رنین \en{$\uparrow$}؛ \en{$\mathrm{BP}\downarrow$}.
  \item \en{$\mathrm{ACTH}\uparrow$} \adarrow\ هایپرپیگمنتیشن پوستی \en{$\blacksquare$}.
  \item \aduncertain{نوشتهٔ کاملاً خط‌خورده در حاشیهٔ چپ این شاخه: ناخوانا}.
  \end{itemize}
\item نارسایی ثانویه \adarrow\ \en{$\mathrm{ACTH}\downarrow$} \adarrow\ گلوکو \en{$\downarrow$} ولی \admark{مینرالو \en{N/L}}:
  \begin{itemize}
  \item \en{$\mathrm{Na}\downarrow$}، \admark{\en{K (N/L)}}.
  \item شایع‌ترین علت \adarrow\ قطع مصرف گلوکوکورتیکوئید خوراکی.
  \end{itemize}
\item \admark{\en{BP} در هر دو نوع نارسایی \en{$\downarrow$} است.}
\item علائم بالینی نارسایی آدرنال \adarrow\ کاهش وزن، ضعف، استفراغ، تمایل شدید به خوردن نمک:
  \begin{itemize}
  \item کمی غیراختصاصی.
  \item تمایل به خوردن نمک: \admark{صرفاً اولیه}.
  \end{itemize}
\item هایپرپیگمانتاسیون:
  \begin{enumerate}
  \item \admark{سطوح اکستانسور}.
  \item چین‌های کف دست.
  \item مخاط دهان.
  \end{enumerate}
\item \en{lab data}:
  \begin{itemize}
  \item \admark{متابولیک اسیدوز خفیف}؛ بدون \en{AG}؛ \admark{صرفاً اولیه}.
  \item \en{$\mathrm{BUN}\uparrow\uparrow$}، \en{$\mathrm{Ca}\uparrow$}، \en{$\mathrm{Hb}\downarrow$}، \en{$\mathrm{lymph}\uparrow$}، \en{$\mathrm{Eosi}\uparrow$}.
  \item \admark{\en{$\mathrm{BS}\downarrow$} (خصوصاً در اطفال)}.
  \item \aduncertain{عبارت خط‌خورده کنار اسیدوز: ناخوانا}.
  \end{itemize}
\item در نارسایی ثانویه، به دلیل \en{$\mathrm{ACTH}\downarrow$} ممکن است علائم \admark{هایپوتیروئیدی، هیپوگنادیسم و کمبود هورمون رشد} وجود داشته باشد.
\end{itemize}
\subsection*{نارسایی حاد آدرنال}
\begin{itemize}
\item نمونهٔ خون جهت \en{ACTH}، کورتیزول، آلدوسترون، رنین؛ سپس بلافاصله \admark{\en{100 mg} هیدروکورتیزون \en{IV}} به همراه \admark{\en{N/S}}.
\item در سپسیس نارسایی حاد آدرنال دیده می‌شود.
\end{itemize}
\adfooter{۱}

\adpage{۲}
\subsection*{نارسایی مزمن آدرنال}
\begin{itemize}
\item \admark{\en{ACTH} را اندازه‌گیری می‌کنیم}.
\item بررسی \adarrow\ \admark{کورتیزول پایهٔ صبحگاهی پلاسما}، \admark{تست کوزین‌تروپین یک‌ساعته}:
\end{itemize}
\begin{latin}
\noindent 0.25 mg ACTH IV then Cortisol level at 0, 30, 60.
\end{latin}
\begin{itemize}
\item کورتیزول پلاسما:
  \begin{itemize}
  \item در حالت طبیعی، در هر زمان: \en{$>20\,\mu\mathrm{g/dl}$}.
  \item نارسایی مزمن آدرنال: صبح \en{$<5\,\mu\mathrm{g/dl}$}؛ تست \en{$<18\,\mu\mathrm{g/dl}$}.
  \item صبح \aduncertain{\en{10--18}} و تحریک‌شده \en{$<18$} \adarrow\ اختلال در ذخیرهٔ آدرنال.
  \item \en{In stress} \adarrow\ کورتیزول می‌دهیم.
  \end{itemize}
\end{itemize}
\subsection*{سندرم‌های همراه با آدیسون}
\admark{در صورت شک، محور آدرنال را بررسی می‌کنیم.}
\begin{latin}
\noindent Polyglandular Autoimmune
\end{latin}
\begin{itemize}
\item \en{PGA-Type 1} \adarrow\ هایپوتیروئیدی، نارسایی آدرنال، \admark{کاندیدیازیس جلدی مخاطی}؛ \admark{در اطفال} ایجاد می‌شود.
\item \en{PGA-Type 2} (\admark{سندرم اشمیت}) \adarrow\ آدیسون، \en{DM} وابسته به انسولین، بیماری تیروئید:
  \begin{itemize}
  \item گریوز، هاشیموتو.
  \end{itemize}
\end{itemize}
\subsection*{نارسایی ثانویهٔ آدرنال}
\begin{itemize}
\item اسیدوز متابولیک \en{$\times$}؛ هایپرکالمی \en{$\times$}؛ میل به نمک \en{$\times$}.
\item هایپوتیروئیدی، کاهش هورمون رشد، هایپوگنادیسم \adarrow\ صرفاً در نارسایی ثانویه.
\item درمان نگهدارندهٔ بیماری آدیسون (نارسایی اولیه) \adarrow\ گلوکوکورتیکوئید + فلودروکورتیزون \en{life long}؛ مینرالو.
\item تست کوزین‌تروپین علت اولیه و ثانویه را افتراق نمی‌دهد.
\item نارسایی اولیه \adarrow\ \en{$\mathrm{ACTH}>20\,\mathrm{pg/ml}$}، فعالیت رنین ایستاده \en{$>3\,\mathrm{ng/ml/h}$}.
\item نارسایی ثانویه \adarrow\ \en{$\mathrm{ACTH}<20\,\mathrm{pg/ml}$}، فعالیت رنین ایستاده \en{$<3\,\mathrm{ng/ml/h}$}.
\end{itemize}
\adfooter{۲}

\adpage{۳}
\adprinted
\begin{itemize}
\item اصول درمان نارسایی آدرنال \adarrow\ باید حداقل دوز هیدروکورتیزون داده شود.
\item مینرالو با دوزی که فعالیت رنین پلاسما (\en{PRA}) در محدودهٔ \en{$1\text{--}3\,\mathrm{ng/ml/h}$} باشد.
\item کریز حاد آدرنال = نارسایی حاد آدرنال \adarrow\ \admark{هایپوتنشن شدید، افت هوشیاری، کوما، هایپوگلیسمی}.
\end{itemize}
\subsection*{هایپرپلازی مادرزادی آدرنال (\en{CAH})}
\begin{itemize}
\item \admark{شایع‌ترین نوع: کمبود ۲۱-هیدروکسیلاز} (۹۵\% موارد).
\item \admark{کمبود کورتیزول و آلدوسترون}.
\item \en{$\mathrm{ACTH}\uparrow$} \adarrow\ ۱۷-هیدروکسی‌پروژسترون \en{$\uparrow$}، پروژسترون \en{$\uparrow$}.
\item آندروستندیون \en{$\uparrow$}، \en{$\mathrm{DHEA}\uparrow$} \adarrow\ \admark{تستوسترون \en{$\uparrow$}}.
\end{itemize}
\subsection*{کمبود ۲۱-هیدروکسیلاز}
\begin{itemize}
\item \admark{\en{$\mathrm{K}\uparrow$}}.
\item نوع کلاسیک \adarrow\ تشخیص در تولد یا کودکی.
\item \en{$2/3$} موارد کمبود مینرالو \en{$\checkmark$}.
\item \en{$1/3$} موارد کمبود مینرالو \en{$\times$}، علائم ویریلیزاسیون \en{$\checkmark$}.
\item \admark{تست تشخیصی} \adarrow\ \admark{۱۷-هیدروکسی‌پروژسترون \en{$>200$}}.
\item نوع دیررس \adarrow\ شایع‌ترین بیماری \en{AR} در انسان.
\item حوالی بلوغ \adarrow\ ویریلیزیشن.
\item هیرسوتیسم، آمنوره، الیگومنوره، ناباروری، اختلال قاعدگی.
\item تشخیص \adarrow\ \en{17-OH} پرو \en{$>1500$}، طی \en{30 min} بعد از تجویز \en{0.25} میلی‌گرم \en{ACTH}.
\item درمان کلیهٔ موارد \en{CAH} \adarrow\ \admark{گلوکومینرالوکورتیکوئید}.
\end{itemize}
\subsection*{کمبود ۱۱-هیدروکسیلاز}
\begin{itemize}
\item ۵\% موارد \en{CAH}.
\item \admark{کورتیزول \en{$\downarrow$}، دئوکسی‌کورتیکوسترون (پیش‌ساز آلدوسترون)}.
\item \en{HTN}، \en{$\mathrm{K}\downarrow$}، \admark{ویریلیزاسیون}.
\item نوع دیررس \adarrow\ صرفاً ویریلیزاسیون در حوالی بلوغ.
\end{itemize}
\adfooter{۳}

\adpage{۴}
\begin{itemize}
\item تشخیص \adarrow\ \admark{افزایش ۱۱-دئوکسی‌کورتیزول پلاسما} در حالت پایه و پس از تحریک با \en{ACTH}.
\end{itemize}
\subsection*{کمبود آنزیم‌ها}
\begin{center}
\begin{tabular}{|r|c|c|c|}
\hline
کمبود آنزیم‌ها & \en{BP} & \en{K} & ویریلیزیشن \\
\hline
۲۱-هیدروکسیلاز & \admark{\en{$\downarrow$}} & \en{$\uparrow$} & \en{$\checkmark$} \\
۱۱-هیدروکسیلاز & \en{$\uparrow$} & \en{$\downarrow$} & \en{$\checkmark$} \\
۱۷-هیدروکسیلاز & \en{$\uparrow$} & \en{$\downarrow$} & \admark{\en{$\times$}} \\
\hline
\end{tabular}
\end{center}
\subsection*{افزایش سطح کورتیزول}
\begin{itemize}
\item فیزیولوژیک \adarrow\ استرس، \admark{سه‌ماههٔ سوم حاملگی}، \aduncertain{ورزش سنگین و منظم}.
\item پاتولوژیک \adarrow\ سندرم کوشینگ، محرومیت از الکل یا نارکوتیک‌ها.
\item اختلالات روانی: افسردگی، مانیک، الکلیسم، آنورکسیا نرووزا.
\end{itemize}
\subsection*{سندرم کوشینگ}
\begin{itemize}
\item چاقی مرکزی (۹۵\% بیماران)، لاغرشدن اندام‌ها، قرمزشدن صورت (\en{Facial Pletora})، چربی سوپراکلاویکولار.
\item \en{Moon Face}، بوفالو هامپ، ضعف عضلات پروگزیمال.
\item بی‌خوابی، برانگیختگی در شب و غروب، تغییرات خلق.
\item \admark{قاعدگی نامنظم (۸۰\%)}، کاهش لیبیدو، نازک‌شدن پوست (۸۰\%)؛ در پشت دست‌ها.
\item هیرسوتیسم (۷۰\%)، آکنه (۵۰\%)، استریا بنفش (۹۵\%)؛ شکم و ران.
\item \en{lab data} \adarrow\ \admark{\en{$\mathrm{Alk\!\!\!-\!p}\uparrow$}}، \admark{نوتروفیلوز}، \admark{ترومبوسیتوز}، هایپرلیپیدمی، هایپر \en{TG}، \en{DM}.
\item هایپرکورتیزولیسم ناشی از ترشح \en{ACTH} اکتوپیک \adarrow\ \admark{\en{$\mathrm{K}\downarrow$}، آلکالوز}.
\end{itemize}
\subsection*{تست کورتیزول آزاد ادرار (\en{UFC})}
\begin{itemize}
\item در ادرار ۲۴ساعته.
\item \admark{حساسیت ۹۰\%} \adarrow\ \admark{\en{$\mathrm{UFC}>50\,\mu\mathrm{g}/24\mathrm{h}$}}.
\end{itemize}
\adfooter{۴}

\adpage{۵}
\adprinted
\subsection*{\en{Low dose Dexa suppression test (1 mg Dexa)}}
\begin{itemize}
\item ساعت ۱۱ شب دگزا \adarrow\ ساعت ۸ صبح کورتیزول پلاسما.
\item \en{if $>1.8\,\mu\mathrm{g/dl}$: $+$}.
\item مثبت و منفی زیاد.
\item تست کورتیزول شبانه \adarrow\ اگر کورتیزول شبانهٔ پلاسما بیش از ۵۰\% کورتیزول صبح می‌باشد.
\item تست کورتیزول بزاق شبانه \adarrow\ \admark{حساسیت و اختصاصیت بالا}.
\end{itemize}
\subsection*{پس از تأیید سندرم کوشینگ}
\begin{itemize}
\item اولین قدم \adarrow\ \admark{سنجش \en{ACTH}}.
\item \en{ACTH} \adarrow\ \en{N/L}، \en{$\uparrow$} \adarrow\ تومور هیپوفیز یا سندرم ترشح اکتوپیک \en{ACTH}.
\item \en{$\mathrm{ACTH}\downarrow$} \adarrow\ پاتولوژی در آدرنال \adarrow\ \en{CT} یا \en{MRI}.
\end{itemize}
\subsection*{علل سندرم کوشینگ}
% The ectopic-column downward ACTH arrow is as drawn in the source.
\begin{center}
\begin{tabular}{|r|c|c|c|}
\hline
 & آدرنال & هیپوفیز & اکتوپیک \\
\hline
\en{ACTH} & \en{$\downarrow$} & \en{$\uparrow$} & \en{$\uparrow$} \\
\hline
\en{IPSS} & جایگاهی ندارد & \en{$\mathrm{ACTH}\uparrow$} & \en{$\mathrm{ACTH}\downarrow$} \\
\hline
\en{High dose Dexa (8 mg)} & \en{$\times$}؛ سرکوب نمی‌شود & \admark{\en{$\checkmark$}} & \en{$\times$} \\
\hline
\end{tabular}
\end{center}
\en{IPSS}: بررسی خون سینوس پتروزال.
\subsection*{هایپرآلدوسترونیسم}
\begin{itemize}
\item علائم افزایش مینرالوکورتیکوئیدها \adarrow\ احتباس نمک، \en{HTN}، \en{$\mathrm{K}\downarrow$}، \admark{آلکالوز متابولیک}.
\item اتیولوژی:
  \begin{itemize}
  \item \admark{آدنوم آدرنال (شایع‌ترین)}.
  \item هایپرپلازی دوطرفهٔ آدرنال \adarrow\ کورتکس آدرنال.
  \end{itemize}
\end{itemize}
\adfooter{۵}

\adpage{۶}
\begin{itemize}
\item علائم هایپوکالمی خفیف \adarrow\ \admark{خستگی} (\admark{شایع‌ترین شکایت})، ضعف عضلانی، شب‌ادراری، بی‌حالی، سردرد.
\item علائم هایپوکالمی شدید \adarrow\ پرنوشی، پرادراری، پارستزی، تتانی، فلج متناوب.
\item علائم اصلی هایپرآلدو \adarrow\ به علت هایپوکالمی.
\item آزمایشگاه \adarrow\ \admark{آلدوسترون \en{$\uparrow$}، رنین \en{$\downarrow$}}؛ \admark{\en{$\mathrm{Ald}>15$}}؛ \admark{\en{$\mathrm{Ald}/\mathrm{PRA}>20$}}.
  \begin{itemize}
  \item اصلاح دست‌نوشتهٔ مخرج نسبت: \en{Ren} خط خورده و \en{PRA} جایگزین شده؛ عدد پیش از ۲۰ خط خورده و ناخواناست.
  \end{itemize}
\item آدنوم منفرد آدرنال \adarrow\ آدرنالکتومی یک‌طرفه.
\item هایپرپلازی دوطرفه \adarrow\ دیورتیک نگهدارندهٔ پتاسیم: \admark{اسپیرونولاکتون، اپلرنون}.
\item در هایپرآلدو اولیه \adarrow\ \admark{حتماً رنین \en{$\downarrow$} است}.
\item صرفاً به دلیل \admark{\en{$\mathrm{Ca}\downarrow$} آزاد} \adarrow\ \aduncertain{واژهٔ ناخوانا} و ترمور.
\end{itemize}
\subsection*{تشخیص هایپرآلدوسترونیسم}
\begin{enumerate}
\item \admark{گام اول: اثبات وجود \en{$\mathrm{K}\downarrow+\mathrm{HTN}$}}.
\item باید دیورتیک به مدت ۳ هفته قطع شود.
\item گام بعد \adarrow\ \admark{سنجش \en{$\dfrac{\mathrm{Ald}}{\mathrm{PRA}}$}}.
\item سپس \admark{تست‌های تأییدی}:
  \begin{itemize}
  \item تجویز خوراکی سدیم.
  \item انفوزیون سالین.
  \item سرکوب فلودروکورتیزون.
  \item کاپتوپریل \en{challenge}.
  \end{itemize}
\item گام بعد \adarrow\ \admark{\en{CT} یا \en{MRI} آدرنال}.
\item گام بعد \adarrow\ \admark{نمونه‌گیری وریدی آدرنال}:
  \begin{itemize}
  \item آدرنالکتومی یک‌طرفه.
  \item اسپیرونولاکتون یا اپلرنون.
  \end{itemize}
\end{enumerate}
{\footnotesize حاشیهٔ چپ: \aduncertain{۳}، \en{P} در دایره، \aduncertain{علامت ناخوانا کنار سطر آخر}.\par}
\adfooter{۶}

\adpage{۷}
\adprinted
\subsection*{فئوکروموسایتوم}
\begin{itemize}
\item در ۹۰\% موارد منشأ آن \admark{مدولای} آدرنال است.
\item شایع‌ترین علامت: حملات فشار خون.
\item تریاد: سردرد، تپش قلب، تعریق.
\item اضطراب، \aduncertain{ضعف}، کاهش وزن، تهوع.
\item درد شکم و قفسهٔ سینه.
\item عوامل \en{Flair} \adarrow\ استرس روانی، بیهوشی، افزایش فشار داخل شکم.
\item مصرف غذاهای حاوی \admark{تیرامین}.
\end{itemize}
\subsection*{تشخیص فئوکروموسایتوم}
\begin{itemize}
\item غربالگری: \admark{متانفرین و کاتکول‌آمین‌های ادرار}.
\item بهترین تست برای تأیید یا \en{R/O} \adarrow\ \admark{متانفرین و نورمتانفرین آزاد پلاسما}:
  \begin{itemize}
  \item نورمتانفرین \en{$>0.61\,\mathrm{nmol/L}$} \aduncertain{(رقم آستانه)}.
  \item متانفرین \en{$>0.31\,\mathrm{nmol/L}$}.
  \end{itemize}
\item اگر افزایش مختصر باشد \adarrow\ تست مهار کلونیدین.
\item پس از تأیید تشخیص \adarrow\ \admark{\en{CT} آدرنال}.
\item اگر منفی \adarrow\ اسکن \en{$^{131}\mathrm{I}$-MIBG}، اسکن اکتروتاید، \en{MRI}، \en{PET Scan}.
\item \en{MRI} \adarrow\ \en{$T_2$-weighted: High signal}.
\end{itemize}
\subsection*{همراهی فئو با \en{MEN-2}}
\begin{itemize}
\item \admark{\en{MEN-2B} (سندرم سیپل)}:
  \begin{itemize}
  \item گانگلیونوروم روده‌ای، نوروم مخاطی، \en{Medullary Thyroid carcinoma}.
  \item علائم شبه مارفان، فئو.
  \end{itemize}
\item \admark{\en{MEN-2A}} \adarrow\ \admark{فئو، \en{MTC}، هایپرپارا}.
\item فئو با بیماری‌های زیر همراهی دارد \adarrow\ \admark{\en{NF}}، \admark{\en{TS}}، \admark{سندرم فون هیپل لیندو}.
\item فون هیپل لیندو \adarrow\ \aduncertain{همانژیوبلاستوز مخچه‌ای ـ شبکیه‌ای}.
\end{itemize}
\adfooter{۷}

\adpage{۸}
{\footnotesize حاشیهٔ بالای درمان: \aduncertain{ن در دایره}.\par}
\subsection*{درمان}
\begin{itemize}
\item درمان \adarrow\ جراحی.
\item قبل از جراحی:
  \begin{itemize}
  \item \en{$\alpha$}-بلاکر (فنوکسی‌بنزامین) \adarrow\ از ۱--۲ هفته قبل از جراحی.
  \item \en{$\beta$}-بلاکر \adarrow\ قبل و حین جراحی.
  \end{itemize}
\end{itemize}
\subsection*{\en{Incidentaloma}}
\begin{itemize}
\item سندرم کوشینگ تحت‌بالینی \adarrow\ سنجش \en{ACTH} پلاسما، \en{low dose Dexa suppression test}.
\item اگر \en{HTN} وجود داشت \adarrow\ \admark{پتاسیم سرم}، \admark{\en{$\dfrac{\mathrm{Ald}}{\mathrm{PRA}}$}}، \admark{متانفرین آزاد پلاسما یا ادرار}.
\item اندیکاسیون جراحی \adarrow\ \admark{\en{$>4\,\mathrm{cm}$}}، \admark{وجود فعالیت هورمونی}.
\item اگر اندیکاسیون جراحی نداشته باشد \adarrow\ پیگیری، تصویربرداری و سنجش هورمون.
\end{itemize}
\adfooter{۸}

% Printed Date:/Subject: labels occur on odd source pages; GNOTE and page
% numbers occur on every source page. Blank ruled space contains no text.
% Obliterated and uncertain handwriting is explicitly marked in the text.
\ifnum\AdrenalStandalone=1
  \let\AdrenalEnd\enddocument
\else
  \let\AdrenalEnd\relax
\fi
\AdrenalEnd
